\chapter{Stasheff Data}

\begin{definition}[Operation target degree]
\lean{AInfinityTheory.operationTargetDeg}
\label{AInfinityTheory.operationTargetDeg}
\leanok
\uses{AInfinityTheory.shift_ofInt}
Given input degrees $(d_0,\dots,d_{n-1})$, the output degree of the $n$-ary
operation is
\[
  \sum_i d_i + (2-n),
\]
where the final summand is interpreted using the grading shift
\lean{AInfinityTheory.shift_ofInt}.
\end{definition}

\begin{definition}[Stasheff target degree]
\lean{AInfinityTheory.stasheffTargetDeg}
\label{AInfinityTheory.stasheffTargetDeg}
\leanok
\uses{AInfinityTheory.shift_ofInt}
The common degree of the arity-$n$ Stasheff relation is
\[
  \sum_i d_i + (3-n)
\]
in the grading index.
\end{definition}

\begin{definition}[Composable morphism type]
\lean{AInfinityTheory.composableHomType}
\label{AInfinityTheory.composableHomType}
\leanok
\uses{AInfinityTheory.GradedRModule}
For a string of objects and prescribed degrees,
\lean{AInfinityTheory.composableHomType} is the $i$-th graded hom-module in
that composable string.
\end{definition}

\begin{definition}[Operation target type]
\lean{AInfinityTheory.operationTargetType}
\label{AInfinityTheory.operationTargetType}
\leanok
\uses{AInfinityTheory.operationTargetDeg, AInfinityTheory.composableHomType}
For a composable string of objects and degrees, this is the graded hom-module in
which the corresponding multilinear operation takes values.
\end{definition}

\begin{definition}[Indexed Stasheff term]
\lean{AInfinityTheory.indexedStasheffTerm}
\label{AInfinityTheory.indexedStasheffTerm}
\leanok
\uses{AInfinityTheory.composableHomType, AInfinityTheory.operationTargetType, AInfinityTheory.stasheffTargetDeg}
Fix object-indexed operations $m_n$. For valid indices $(r,s)$,
\lean{AInfinityTheory.indexedStasheffTerm} is the term obtained by first applying the inner
$s$-ary operation in positions $r,\dots,r+s-1$ and then applying the outer
operation to the collapsed string.
\end{definition}

\begin{definition}[Stasheff sign parity]
\lean{AInfinityTheory.stasheffSignParity}
\label{AInfinityTheory.stasheffSignParity}
\leanok
\uses{AInfinityTheory.Grading}
For a valid pair $(r,s)$, this is the parity
\[
  |a_{r+s}| + \cdots + |a_{n-1}| - (n-r-s)
\]
computed in \lean{AInfinityTheory.Parity}.
\end{definition}

\begin{definition}[Stasheff sign]
\lean{AInfinityTheory.stasheffSign}
\label{AInfinityTheory.stasheffSign}
\leanok
\uses{AInfinityTheory.stasheffSignParity}
The integer sign attached to $(r,s)$ is $(-1)$ raised to the parity recorded by
\lean{AInfinityTheory.stasheffSignParity}.
\end{definition}

\begin{definition}[Indexed Stasheff sum]
\lean{AInfinityTheory.indexedStasheffSum}
\label{AInfinityTheory.indexedStasheffSum}
\leanok
\uses{AInfinityTheory.indexedStasheffTerm, AInfinityTheory.stasheffSign, AInfinityTheory.stasheffTargetDeg}
The arity-$n$ Stasheff sum is the signed sum of all valid composites of an inner
operation followed by an outer operation.
\end{definition}

\begin{definition}[Indexed Stasheff identities]
\lean{AInfinityTheory.indexedSatisfiesStasheff}
\label{AInfinityTheory.indexedSatisfiesStasheff}
\leanok
\uses{AInfinityTheory.indexedStasheffSum}
An object-indexed family of multilinear operations satisfies the Stasheff
identities if every indexed Stasheff sum vanishes.
\end{definition}
